
\documentclass[preview]{standalone}
\usepackage{amsmath,amssymb, tikz-cd, quiver, pgfplots, mathrsfs}
\usetikzlibrary{lindenmayersystems, nfold}
\usetikzlibrary{decorations.markings,arrows}

\begin{document}
[![](https://q.uiver.app/quiver.svg)](https://github.com/varkor/quiver)

![](https://q.uiver.app/quiver-blue.svg)Loading diagram...

Lightness:

Sync label/edge colours:

Preset:

LaTeX:(None)Diagram:

# Keyboard shortcuts⌘/

For a guide on using keyboard shortcuts in **quiver** see [the tutorial](https://github.com/varkor/quiver/blob/master/tutorial.md).

## General

|   |   |
|---|---|
|Dismiss errors, and panels; Cancel modification or movement; Hide focus point; Deselect, and dequeue cells|esc|
|Import from LaTeX|⌘I|
|Export to LaTeX or Typst|⌘E|

## Navigation

|   |   |
|---|---|
|Pan view|Scroll|
|Enable mouse panning|ctrl,alt|
|Enable touch panning|Long press|
|Enable mouse zooming|⇧|
|Move focus point|←,↑,↓,→|
|Select next queued cell|⇥|
|Select previous queued cell|⇧⇥|
|Select / deselect object|S|
|Select cells|;|
|Toggle cell selection|'|

## Modification

|   |   |
|---|---|
|Focus / defocus label input|↵|
|Create object, and connect to selection||
|Move selected objects|B|
|Change source|,|
|Change target|.|
|Create arrows from selection|/|
|Copy|⌘C|
|Cut|⌘X|
|Paste|⌘V|

## Styling

|   |   |
|---|---|
|Reverse arrows|R|
|Flip arrows|E|
|Flip labels|F|
|Left-align labels|V|
|Centre-align labels|C|
|Over-align labels|X|
|Modify label position|I|
|Modify offset|O|
|Modify curve|K|
|Modify radius|N|
|Modify length|L (hold ⇧ to shorten symmetrically)|
|Modify level|M|
|Modify style|D|
|Display as arrow|A|
|Display as adjunction|J|
|Display as pullback/pushout|P (press again to switch corner style)|
|Modify label colour|Y|
|Modify arrow colour|U|

## Toolbar

|   |   |
|---|---|
|Save diagram in URL|⌘S|
|Undo|⌘Z|
|Redo|⇧⌘Z|
|Select all|⌘A|
|Select connected components|⇧⌘C|
|Deselect all|⇧⌘A|
|Delete|⌫,del|
|Centre view|G|
|Zoom out|⌘-|
|Zoom in|⌘=|
|Toggle grid|H|
|Toggle hints|⇧⌘H|

## Export

|   |   |
|---|---|
|Toggle diagram centring|C|
|Toggle ampersand replacement|A|
|Toggle cramped spacing|R|
|Toggle standalone|S|
|Toggle fixed size|F|

# About

**quiver** is a modern, graphical editor for commutative and pasting diagrams, capable of rendering high-quality diagrams for screen viewing, and exporting to LaTeX via `tikz-cd` and Typst via `fletcher`.

Creating and modifying diagrams with **quiver** is orders of magnitude faster than writing the equivalent LaTeX or Typst by hand and, with a little experience, competes with pen-and-paper. To learn how to use **quiver** efficiently, see [the tutorial](https://github.com/varkor/quiver/blob/master/tutorial.md).

The editor is open source and may be found [on GitHub](https://github.com/varkor/quiver). If you would like to request a feature, or want to report an issue, you can [do so here](https://github.com/varkor/quiver/issues).

You can follow **quiver** on [Mastodon](https://mathstodon.xyz/@quiver) or [Twitter](https://twitter.com/q_uiver_app) for updates on new features.

## Thanks to

- [S. C. Steenkamp](https://www.cl.cam.ac.uk/~scs62/), for helpful discussions regarding the aesthetic rendering of arrows.

- [AndréC](https://tex.stackexchange.com/users/138900/andr%c3%a9c), for the custom TikZ style for curves of a fixed height.

- [Andrew Stacey](https://tex.stackexchange.com/users/86/andrew-stacey), for the custom TikZ style for shortened curves.

- [Théophile Cailliau](https://github.com/tjbcg), for implementing Typst support.

- [Nathan Corbyn](https://github.com/doctorn), for adding the ability to export embeddable diagrams to HTML.

- [Paolo Brasolin](https://github.com/paolobrasolin), for adding offline support.

- [Carl Davidson](https://github.com/davidson16807), for discussing and prototyping loop rendering.

- Everyone who has improved **quiver** by submitting pull requests, reporting issues or suggesting improvements.

Created by [varkor](https://github.com/varkor).

# Welcome

**quiver** is a modern, graphical editor for commutative and pasting diagrams, capable of rendering high-quality diagrams for screen viewing, and exporting to LaTeX via `tikz-cd` and Typst via `fletcher`.

**quiver** is intended to be intuitive to use and easy to pick up. Here are a few tips to help you get started:

- Click and drag to create new arrows: the source and target objects will be created automatically.
- Double-click to create a new object.
- Edit labels with the input bar at the bottom of the screen.
- Click and drag the empty space around a object to move it around.
- Hold Shift (⇧) to select multiple cells to edit them simultaneously.

For a detailed guide to using **quiver**, see [the tutorial](https://github.com/varkor/quiver/blob/master/tutorial.md).

Remember to include `\usepackage{quiver}` in your LaTeX preamble. You can install the package using [CTAN](https://ctan.org/pkg/quiver), or [open `quiver.sty` in a new tab](https://raw.githubusercontent.com/varkor/quiver/master/package/quiver.sty) to copy-and-paste.updated on 2025-07-05

% https://q.uiver.app/#q=WzAsMyxbMCwwLCJwXnstMX0oVSkiXSxbMSwwLCJVXFx0aW1lcyBJIl0sWzEsMSwiVSJdLFswLDIsInAiLDJdLFsxLDIsIlxccGlfVSJdLFswLDEsImgiXV0= \[\begin{tikzcd} {p^{-1}(U)} & {U\times I} \\ & U \arrow["h", from=1-1, to=1-2] \arrow["p"', from=1-1, to=2-2] \arrow["{\pi_U}", from=1-2, to=2-2] \end{tikzcd}\]
\end{document}

